\begin{table}[htbp]\centering
\caption{Additional grid coverage}\label{tab:review_appendix_design}
\begin{minipage}{.95\textwidth}\fontsize{8}{10}\selectfont\setstretch{1}\setlength{\tabcolsep}{2pt}\renewcommand{\arraystretch}{1.08}
\par\noindent\textbf{Full supporting domains}\par\smallskip
\begin{tabular*}{\linewidth}{@{\extracolsep{\fill}}lrrrr@{}}\toprule
ID & State jobs & Records & P3 feasible designs & P4 feasible designs \\
\midrule
E4 & 90 & 360 & 52 & 52 \\
E5 & 93 & 744 & 93 & -- \\
E5C & 54 & 432 & 54 & -- \\
E7 & 75 & 300 & 34 & 34 \\
\bottomrule\end{tabular*}\par\medskip
\par\noindent\fontsize{8}{9.5}\selectfont E4 crosses two states, five floors and allowances 0/2,000/5,000 CNY per reference departure. E5 includes 13 numerical/preference/authority configurations and a load--cost grid rebased to each load	extquotesingle s reference revenue. E5C holds unallocated model cost at the central 0.88 calibration and supplies the main structural ranges. Both grids cross three reference loads, two ceilings and three variable shares. E7 crosses polarization and road-saving scales 0/0.25/0.50/1.00/1.50. Counts pool feasible saved designs across objectives; failure to find a feasible design does not prove infeasibility.
\end{minipage}\end{table}
